\textbf{Phần 1. Trắc nghiệm lựa chọn}
\setcounterpageref{ex}{0}

\begin{ex} %[4D10V58]
Trong hai phân số có cùng mẫu số dương, phân số nào lớn hơn?
\choice
{Phân số có tử số bé hơn}
{\True Phân số có tử số lớn hơn}
{Hai phân số luôn bằng nhau}
{Không so sánh được}
\loigiai{
Cùng mẫu số: tử lớn hơn thì phân số lớn hơn.\\
\textbf{Đáp án của câu hỏi là B}
}
\end{ex}

\begin{ex} %[4D10V58]
Điền dấu thích hợp vào chỗ chấm: $\dfrac{3}{7} \dots \dfrac{5}{7}$.
\choice
{\True $<$}
{$>$}
{$=$}
{Không so sánh được}
\loigiai{
$3 < 5 \Rightarrow \dfrac{3}{7} < \dfrac{5}{7}$.\\
\textbf{Đáp án của câu hỏi là A}
}
\end{ex}

\begin{ex} %[4D10V58]
Trong hai phân số có cùng tử số dương, phân số nào lớn hơn?
\choice
{Phân số có mẫu số lớn hơn}
{Hai phân số luôn bằng nhau}
{\True Phân số có mẫu số bé hơn}
{Không so sánh được}
\loigiai{
Cùng tử số: mẫu bé hơn thì phân số lớn hơn.\\
\textbf{Đáp án của câu hỏi là C}
}
\end{ex}

\begin{ex} %[4D10V58]
Điền dấu thích hợp vào chỗ chấm: $\dfrac{4}{5} \dots \dfrac{4}{9}$.
\choice
{$<$}
{$=$}
{Không so sánh được}
{\True $>$}
\loigiai{
$5 < 9 \Rightarrow \dfrac{4}{5} > \dfrac{4}{9}$.\\
\textbf{Đáp án của câu hỏi là D}
}
\end{ex}

\begin{ex} %[4D10V58]
So sánh hai phân số $\dfrac{2}{3}$ và $\dfrac{3}{4}$:
\choice
{$\dfrac{2}{3} > \dfrac{3}{4}$}
{\True $\dfrac{2}{3} < \dfrac{3}{4}$}
{$\dfrac{2}{3} = \dfrac{3}{4}$}
{Không so sánh được}
\loigiai{
$\dfrac{2}{3} = \dfrac{8}{12} < \dfrac{9}{12} = \dfrac{3}{4}$.\\
\textbf{Đáp án của câu hỏi là B}
}
\end{ex}

\begin{ex} %[4D10V58]
Phân số nào sau đây lớn hơn 1?
\choice
{\True $\dfrac{8}{5}$}
{$\dfrac{5}{6}$}
{$\dfrac{7}{7}$}
{$\dfrac{3}{4}$}
\loigiai{
$\dfrac{8}{5} > 1$ (tử số lớn hơn mẫu số).\\
\textbf{Đáp án của câu hỏi là A}
}
\end{ex}

\begin{ex} %[4D10V58]
Sắp xếp các phân số $\dfrac{1}{2}, \dfrac{1}{3}, \dfrac{1}{5}, \dfrac{1}{4}$ theo thứ tự từ bé đến lớn:
\choice
{$\dfrac{1}{2} < \dfrac{1}{3} < \dfrac{1}{4} < \dfrac{1}{5}$}
{$\dfrac{1}{5} < \dfrac{1}{3} < \dfrac{1}{4} < \dfrac{1}{2}$}
{\True $\dfrac{1}{5} < \dfrac{1}{4} < \dfrac{1}{3} < \dfrac{1}{2}$}
{$\dfrac{1}{2} < \dfrac{1}{4} < \dfrac{1}{3} < \dfrac{1}{5}$}
\loigiai{
Cùng tử số: mẫu càng lớn phân số càng bé $\Rightarrow \dfrac{1}{5} < \dfrac{1}{4} < \dfrac{1}{3} < \dfrac{1}{2}$.\\
\textbf{Đáp án của câu hỏi là C}
}
\end{ex}

\begin{ex} %[4D10V58]
So sánh hai phân số $\dfrac{7}{8}$ và $\dfrac{8}{9}$ bằng cách so sánh phần bù tới 1:
\choice
{$\dfrac{7}{8} > \dfrac{8}{9}$}
{$\dfrac{7}{8} = \dfrac{8}{9}$}
{Không so sánh được}
{\True $\dfrac{7}{8} < \dfrac{8}{9}$}
\loigiai{
Phần bù: $1 - \dfrac{7}{8} = \dfrac{1}{8} > \dfrac{1}{9} = 1 - \dfrac{8}{9} \Rightarrow \dfrac{7}{8} < \dfrac{8}{9}$.\\
\textbf{Đáp án của câu hỏi là D}
}
\end{ex}

\begin{ex} %[4D10V58]
Bạn An ăn $\dfrac{1}{3}$ cái bánh, bạn Bình ăn $\dfrac{2}{5}$ cái bánh. Bạn nào ăn nhiều bánh hơn?
\choice
{Bạn An}
{\True Bạn Bình}
{Hai bạn ăn bằng nhau}
{Không so sánh được}
\loigiai{
$\dfrac{1}{3} = \dfrac{5}{15} < \dfrac{6}{15} = \dfrac{2}{5}$ nên bạn Bình ăn nhiều hơn.\\
\textbf{Đáp án của câu hỏi là B}
}
\end{ex}

\begin{ex} %[4D10V58]
Phân số nào sau đây nằm giữa $\dfrac{1}{3}$ và $\dfrac{2}{3}$?
\choice
{\True $\dfrac{1}{2}$}
{$\dfrac{1}{6}$}
{$\dfrac{5}{6}$}
{$\dfrac{3}{4}$}
\loigiai{
$\dfrac{1}{3} = \dfrac{2}{6} < \dfrac{1}{2} = \dfrac{3}{6} < \dfrac{2}{3} = \dfrac{4}{6}$.\\
\textbf{Đáp án của câu hỏi là A}
}
\end{ex}

\begin{ex} %[4D10V58]
Tìm phân số lớn nhất trong các phân số: $\dfrac{3}{4}, \dfrac{5}{6}, \dfrac{2}{3}, \dfrac{7}{12}$.
\choice
{$\dfrac{3}{4}$}
{$\dfrac{2}{3}$}
{\True $\dfrac{5}{6}$}
{$\dfrac{7}{12}$}
\loigiai{
MSC là 12: $\dfrac{9}{12}, \dfrac{10}{12}, \dfrac{8}{12}, \dfrac{7}{12}$. Lớn nhất là $\dfrac{10}{12} = \dfrac{5}{6}$.\\
\textbf{Đáp án của câu hỏi là C}
}
\end{ex}

\begin{ex} %[4D10V58]
So sánh hai phân số: $A = \dfrac{2025}{2026}$ và $B = \dfrac{2026}{2027}$.
\choice
{$A > B$}
{$A = B$}
{Không so sánh được}
{\True $A < B$}
\loigiai{
Phần bù: $1 - A = \dfrac{1}{2026} > \dfrac{1}{2027} = 1 - B \Rightarrow A < B$.\\
\textbf{Đáp án của câu hỏi là D}
}
\end{ex}

\textbf{Phần 2. Trắc nghiệm đúng sai}
\setcounterpageref{ex}{0}

\begin{ex} %[4D10V58]
Xét tính đúng sai của các phép so sánh phân số:
\choiceTF
{\True $\dfrac{5}{8} < \dfrac{7}{8}$.}
{\True $\dfrac{3}{4} > \dfrac{3}{7}$.}
{$\dfrac{5}{4} < 1$.}
{\True $\dfrac{8}{8} = 1$.}
\loigiai{
\begin{itemize}
    \item $5/8 < 7/8$: Đúng.
    \item $3/4 > 3/7$: Đúng.
    \item $5/4 > 1$ (không bé hơn 1). Sai.
    \item $8/8 = 1$: Đúng.
\end{itemize}
\begin{itemize}
    \item a) \textbf{Đúng}.
    \item b) \textbf{Đúng}.
    \item c) \textbf{Sai}.
    \item d) \textbf{Đúng}.
\end{itemize}
}
\end{ex}

\begin{ex} %[4D10V58]
Cho 3 phân số: $\dfrac{3}{5}, \dfrac{2}{3}, \dfrac{7}{10}$. Xét tính đúng sai:
\choiceTF
{\True Quy đồng với mẫu số chung 30 được $\dfrac{18}{30}, \dfrac{20}{30}, \dfrac{21}{30}$.}
{\True Phân số lớn nhất là $\dfrac{7}{10}$.}
{Phân số bé nhất là $\dfrac{2}{3}$.}
{\True Sắp xếp tăng dần: $\dfrac{3}{5} < \dfrac{2}{3} < \dfrac{7}{10}$.}
\loigiai{
\begin{itemize}
    \item $3/5 = 18/30; 2/3 = 20/30; 7/10 = 21/30$. Đúng.
    \item $7/10$ lớn nhất: Đúng.
    \item Bé nhất là $3/5$ (không phải $2/3$). Sai.
    \item Thứ tự tăng dần: $3/5 < 2/3 < 7/10$. Đúng.
\end{itemize}
\begin{itemize}
    \item a) \textbf{Đúng}.
    \item b) \textbf{Đúng}.
    \item c) \textbf{Sai}.
    \item d) \textbf{Đúng}.
\end{itemize}
}
\end{ex}

\textbf{Phần 3. Trả lời ngắn}
\setcounterpageref{ex}{0}

\begin{ex} %[4D10V58]
Trong hai phân số $\dfrac{4}{9}$ và $\dfrac{7}{9}$, phân số nào lớn hơn? (Viết dạng a/b).
\shortans{7/9}
\loigiai{
\textbf{Đáp số:} $7/9$\\
$7/9$.
}
\end{ex}

\begin{ex} %[4D10V58]
Trong hai phân số $\dfrac{5}{6}$ và $\dfrac{5}{8}$, phân số nào lớn hơn? (Viết dạng a/b).
\shortans{5/6}
\loigiai{
\textbf{Đáp số:} $5/6$\\
$5/6$.
}
\end{ex}

\begin{ex} %[4D10V58]
Tìm phân số lớn nhất trong 3 phân số: $\dfrac{1}{2}, \dfrac{2}{3}, \dfrac{3}{4}$ (dạng a/b).
\shortans{3/4}
\loigiai{
\textbf{Đáp số:} $3/4$\\
$3/4$.
}
\end{ex}

\begin{ex} %[4D10V58]
Tìm phân số bé nhất trong 3 phân số: $\dfrac{2}{5}, \dfrac{1}{3}, \dfrac{4}{15}$ (dạng a/b).
\shortans{4/15}
\loigiai{
\textbf{Đáp số:} $4/15$\\
$6/15, 5/15, 4/15 \Rightarrow 4/15$ bé nhất.
}
\end{ex}

\begin{ex} %[4D10V58]
Có bao nhiêu phân số có mẫu số bằng 8 lớn hơn $\dfrac{2}{8}$ và bé hơn $\dfrac{7}{8}$?
\shortans{4}
\loigiai{
\textbf{Đáp số:} $4$\\
$3/8, 4/8, 5/8, 6/8$ (4 phân số).
}
\end{ex}

\begin{ex} %[4D10V58]
Tìm chữ số tự nhiên $x$ sao cho $\dfrac{3}{5} < \dfrac{x}{10} < \dfrac{4}{5}$.
\shortans{7}
\loigiai{
\textbf{Đáp số:} $7$\\
$6/10 < x/10 < 8/10 \Rightarrow x = 7$.
}
\end{ex}
